\section{End-Term Examination --- Set D}
\label{sec:endterm_set_d}

\LPUPaperHeader{End-Term Examination}{D}{3 Hours}{60 Marks}{CO1 to CO6 (Units 1 to 6 Comprehensive)}

\begin{center}
\sffamily\textbf{\color{AlertRed}Structure of Question Paper:}\\[1mm]
\textbf{Part A (Objective):} 30 MCQs strictly from \textbf{Units 4, 5, and 6}. Total weight: \textbf{20 Marks}. All questions compulsory.\\
\textbf{Part B (Subjective):} 6 Questions from \textbf{all units} (10 Marks each). \textbf{Attempt any 4 questions} ($4 \times 10 = \mathbf{40\text{ Marks}}$).\\
\textbf{Total Maximum Marks: 60 Marks.}
\end{center}

\vspace{3mm}

% ------------------------------------------------------------------------------
% PART A: 30 MCQS (UNITS 4, 5, 6)
% ------------------------------------------------------------------------------
\subsection*{PART A: Objective Section (30 MCQs --- Units 4, 5, 6) [20 Marks]}

\mcqitem{1}{If $f(x,y) = \ln(x^2 + y^2)$, then $f_{xx} + f_{yy}$ is equal to:}{$0$}{$1$}{$\frac{2}{x^2+y^2}$}{$\frac{4}{x^2+y^2}$}{1}

\mcqitem{2}{The total differential $dz$ of $z = x^y$ is:}{$y x^{y-1} dx + x^y \ln x \, dy$}{$x^y dx + x^y dy$}{$y x^{y-1} dx$}{$x^y \ln x \, dy$}{1}

\mcqitem{3}{The maximum value of the function $f(x,y) = 2 - x^2 - y^2$ is:}{$2$}{$0$}{$1$}{$-2$}{1}

\mcqitem{4}{The divergence of the vector field $\vec{F} = (x^2 - y^2)\hat{i} + 2xy\hat{j}$ is:}{$4x$}{$2x$}{$0$}{$2y$}{1}

\mcqitem{5}{The curl of the vector field $\vec{F} = (x^2 - y^2)\hat{i} + 2xy\hat{j}$ is:}{$4y\hat{k}$}{$0$}{$2y\hat{k}$}{$-4y\hat{k}$}{1}

\mcqitem{6}{The value of $\iint_R xy \, dx \, dy$ over the rectangle $0 \le x \le 2, 0 \le y \le 4$ is:}{$16$}{$8$}{$32$}{$4$}{1}

\mcqitem{7}{The value of $\int_0^\infty \int_0^\infty e^{-(x+y)} \, dx \, dy$ is:}{$1$}{$0$}{$\infty$}{$1/2$}{1}

\mcqitem{8}{The value of the line integral $\int_C (x dx + y dy)$ along any path from $(0,0)$ to $(1,1)$ is:}{$1$}{$0$}{$1/2$}{$2$}{1}

\mcqitem{9}{If $\vec{F}$ is conservative, then the circulation $\oint_C \vec{F} \cdot d\vec{r}$ around any closed path is:}{$0$}{$2\pi$}{$\pi$}{Area enclosed}{1}

\mcqitem{10}{The volume of a sphere of radius $R$ is calculated by evaluating $\iiint_V dV$. Its value is:}{$\frac{4}{3}\pi R^3$}{$4\pi R^2$}{$\frac{2}{3}\pi R^3$}{$\pi R^3$}{1}

\mcqitem{11}{The value of $\lim_{(x,y)\to(0,0)} \frac{x^2 - y^2}{x^2 + y^2}$ along the line $y = 0$ is:}{$1$}{$-1$}{$0$}{Undefined}{1}

\mcqitem{12}{The directional derivative is zero in the direction:}{Tangent to the level curve}{Parallel to $\nabla \phi$}{Anti-parallel to $\nabla \phi$}{Normal to the surface}{1}

\mcqitem{13}{The normal vector to the surface $x^2 + y^2 - z = 0$ at $(1, 2, 5)$ is:}{$2\hat{i} + 4\hat{j} - \hat{k}$}{$\hat{i} + 2\hat{j} - \hat{k}$}{$2\hat{i} + 4\hat{j} + \hat{k}$}{$\hat{i} + \hat{j} + \hat{k}$}{1}

\mcqitem{14}{The value of $\nabla \cdot (\phi \vec{A})$ is:}{$\phi(\nabla \cdot \vec{A}) + \vec{A} \cdot \nabla\phi$}{$\phi(\nabla \cdot \vec{A})$}{$\vec{A} \cdot \nabla\phi$}{$\nabla\phi \times \vec{A}$}{1}

\mcqitem{15}{The value of $\iint_S \vec{F} \cdot \hat{n} \, dS$ where $\vec{F} = \text{curl} \vec{A}$ over any closed surface is:}{$0$}{$1$}{Volume of $S$}{Area of $S$}{1}

\mcqitem{16}{The Jacobian $\frac{\partial(x,y,z)}{\partial(\rho,\theta,\phi)}$ for spherical polar coordinates is:}{$\rho^2 \sin\phi$}{$\rho \sin\phi$}{$\rho^2 \cos\phi$}{$\rho^2$}{1}

\mcqitem{17}{The value of $\int_0^1 \int_0^1 \int_0^1 dx \, dy \, dz$ is:}{$1$}{$0$}{$3$}{$1/6$}{1}

\mcqitem{18}{The stationary points of $f(x,y) = x^3 + y^3 - 3xy$ are:}{$(0,0)$ and $(1,1)$}{$(0,0)$ only}{$(1,1)$ only}{$(1,-1)$}{1}

\mcqitem{19}{If $u = x^2$ and $v = y^2$, then $\frac{\partial(u,v)}{\partial(x,y)}$ is:}{$4xy$}{$2xy$}{$x+y$}{$2(x+y)$}{1}

\mcqitem{20}{The line integral $\int_C \vec{F} \cdot d\vec{r}$ represents:}{Work done by force $\vec{F}$}{Flux of $\vec{F}$}{Area}{Volume}{1}

\mcqitem{21}{The value of $\nabla(x+y+z)$ is:}{$\hat{i} + \hat{j} + \hat{k}$}{$3$}{$\vec{0}$}{$x\hat{i} + y\hat{j} + z\hat{k}$}{1}

\mcqitem{22}{The area of the circle $r = 2a\cos\theta$ is:}{$\pi a^2$}{$2\pi a^2$}{$4\pi a^2$}{$\frac{\pi a^2}{2}$}{1}

\mcqitem{23}{The vector field $\vec{F} = y\hat{i} + x\hat{j}$ is:}{Irrotational and Conservative}{Solenoidal only}{Neither irrotational nor solenoidal}{Rotational}{1}

\mcqitem{24}{The value of $\iint_S \vec{r} \cdot \hat{n} \, dS$ over a cube of side $L$ is:}{$3L^3$}{$L^3$}{$6L^2$}{$0$}{1}

\mcqitem{25}{Euler's theorem states that for a homogeneous function of degree $n$:}{$x f_x + y f_y = n f$}{$x f_x - y f_y = n f$}{$x^2 f_{xx} = n f$}{$f_x + f_y = n$}{1}

\mcqitem{26}{The value of $\int_0^1 \int_0^x dy \, dx$ is:}{$1/2$}{$1$}{$1/3$}{$2$}{1}

\mcqitem{27}{The value of $\nabla \cdot (\nabla \phi)$ is called the:}{Laplacian of $\phi$}{Curl of $\phi$}{Gradient of $\phi$}{Divergence of $\phi$}{1}

\mcqitem{28}{If $\vec{F} = \nabla \phi$, then $\nabla \times \vec{F}$ is:}{$\vec{0}$}{$1$}{$\nabla^2 \phi$}{$\vec{F}$}{1}

\mcqitem{29}{The transformation $x = r\cos\theta, y = r\sin\theta$ maps the unit disk to:}{$0 \le r \le 1, 0 \le \theta \le 2\pi$}{$0 \le r \le 1, 0 \le \theta \le \pi$}{$-1 \le r \le 1, 0 \le \theta \le 2\pi$}{$r = 1$}{1}

\mcqitem{30}{The value of $\lim_{(x,y)\to(0,0)} \frac{\sin(xy)}{xy}$ is:}{$1$}{$0$}{Undefined}{$\infty$}{1}

\vspace{5mm}
\hrule
\vspace{5mm}

% ------------------------------------------------------------------------------
% PART B: 6 SUBJECTIVE QUESTIONS (ATTEMPT ANY 4)
% ------------------------------------------------------------------------------
\subsection*{PART B: Descriptive Section (Attempt Any 4 of 6) [40 Marks]}

\questionitem{31 (Unit 1)}{
    Investigate for what values of $\lambda$ and $\mu$ the non-homogeneous linear system:
    \[
    \begin{aligned}
    2x + 3y + 5z &= 9 \\
    7x + 3y - 2z &= 8 \\
    2x + 3y + \lambda z &= \mu
    \end{aligned}
    \]
    has: \textbf{(i)} a unique solution, \textbf{(ii)} no solution, and \textbf{(iii)} an infinite number of solutions.
}{10}{CO1}{L3 Apply}

\questionitem{32 (Unit 2)}{
    Solve the \textbf{Cauchy-Euler differential equation}:
    \[
    x^2 \frac{d^2 y}{dx^2} - 3x \frac{dy}{dx} + 4y = 2x^2
    \]
}{10}{CO2}{L3 Apply}

\questionitem{33 (Unit 3)}{
    \textbf{(a)} Test the convergence of the alternating series:
    \[
    \sum_{n=1}^\infty (-1)^{n-1} \frac{n}{2n - 1}
    \]
    \textbf{(b)} Find the interval and radius of convergence of the power series $\sum_{n=1}^\infty \frac{(x+1)^n}{n^2 \cdot 4^n}$.
}{10}{CO3}{L4 Analyze}

\questionitem{34 (Unit 4)}{
    Find the maximum volume of a \textbf{rectangular box open at the top} inscribed in a hemisphere of radius $R$, using multivariable calculus.
}{10}{CO4}{L3 Apply}

\questionitem{35 (Unit 5)}{
    Evaluate the triple integral using spherical coordinates:
    \[
    \iiint_V \frac{dx \, dy \, dz}{\sqrt{1 - x^2 - y^2 - z^2}}
    \]
    over the \textbf{positive octant} of the sphere $x^2 + y^2 + z^2 \le 1$.
}{10}{CO5}{L3 Apply}

\questionitem{36 (Unit 6)}{
    Verify the \textbf{Gauss Divergence Theorem} for the vector field:
    \[
    \vec{F} = x^3\hat{i} + y^3\hat{j} + z^3\hat{k}
    \]
    over the entire solid sphere $V$: $x^2 + y^2 + z^2 \le a^2$.
}{10}{CO6}{L4 Analyze}

\vspace{5mm}
\hrule
\vspace{5mm}

% ------------------------------------------------------------------------------
% COMPLETE STEP-BY-STEP SOLUTIONS
% ------------------------------------------------------------------------------
\subsection*{Solutions \& Marking Scheme (End-Term Set D)}

\subsubsection*{Part A: Answer Key \& Explanations}

\begin{center}
\begin{tabular}{|c|c||c|c||c|c||c|c||c|c||c|c|}
\hline
\textbf{Q} & \textbf{Ans} & \textbf{Q} & \textbf{Ans} & \textbf{Q} & \textbf{Ans} & \textbf{Q} & \textbf{Ans} & \textbf{Q} & \textbf{Ans} & \textbf{Q} & \textbf{Ans} \\
\hline
1 & \textbf{A} & 6 & \textbf{A} & 11 & \textbf{A} & 16 & \textbf{A} & 21 & \textbf{A} & 26 & \textbf{A} \\
2 & \textbf{A} & 7 & \textbf{A} & 12 & \textbf{A} & 17 & \textbf{A} & 22 & \textbf{A} & 27 & \textbf{A} \\
3 & \textbf{A} & 8 & \textbf{A} & 13 & \textbf{A} & 18 & \textbf{A} & 23 & \textbf{A} & 28 & \textbf{A} \\
4 & \textbf{A} & 9 & \textbf{A} & 14 & \textbf{A} & 19 & \textbf{A} & 24 & \textbf{A} & 29 & \textbf{A} \\
5 & \textbf{A} & 10 & \textbf{A} & 15 & \textbf{A} & 20 & \textbf{A} & 25 & \textbf{A} & 30 & \textbf{A} \\
\hline
\end{tabular}
\end{center}

\begin{solutionbox}[Explanations for Part A]
\textbf{Q.1 (Ans: A)} $f_x = \frac{2x}{x^2+y^2} \implies f_{xx} = \frac{2(y^2-x^2)}{(x^2+y^2)^2}$. By symmetry $f_{yy} = \frac{2(x^2-y^2)}{(x^2+y^2)^2}$. Adding gives $0$ (Laplace equation in 2D).\\[1.5mm]
\textbf{Q.4 (Ans: A)} $\nabla \cdot \vec{F} = \frac{\partial}{\partial x}(x^2 - y^2) + \frac{\partial}{\partial y}(2xy) = 2x + 2x = 4x$.\\[1.5mm]
\textbf{Q.8 (Ans: A)} $x dx + y dy = d\left(\frac{x^2+y^2}{2}\right)$. Value is $\frac{1^2+1^2}{2} - 0 = 1$.\\[1.5mm]
\textbf{Q.24 (Ans: A)} $\nabla \cdot \vec{r} = 3$. By Gauss theorem, $\iint_S \vec{r} \cdot \hat{n} dS = \iiint_V 3 dV = 3 \text{Vol}(\text{Cube}) = 3L^3$.
\end{solutionbox}

\subsubsection*{Part B: Detailed Subjective Solutions}

% Solution 31
\begin{solutionbox}[Solution to Question 31 (10 Marks)]
\textbf{Step 1: Augmented Matrix Formulation}
\[
[A|B] = \begin{bmatrix}
2 & 3 & 5 & | & 9 \\
7 & 3 & -2 & | & 8 \\
2 & 3 & \lambda & | & \mu
\end{bmatrix}
\]
Subtract Row 1 from Row 3 ($R_3 \to R_3 - R_1$):
\[
\begin{bmatrix}
2 & 3 & 5 & | & 9 \\
7 & 3 & -2 & | & 8 \\
0 & 0 & \lambda - 5 & | & \mu - 9
\end{bmatrix}
\]
\textbf{Step 2: Checking Rank of Upper $2\times 2$ Submatrix}
\[
\begin{vmatrix} 2 & 3 \\ 7 & 3 \end{vmatrix} = 6 - 21 = -15 \neq 0
\]
Thus, the first two rows are always linearly independent, meaning $\text{rank}(A) \ge 2$ and $\text{rank}[A|B] \ge 2$ for all $\lambda, \mu$.

\textbf{Step 3: Classification of Solutions}
\begin{enumerate}[label=(\roman*)]
    \item \textbf{Unique Solution:}
    Requires $\text{rank}(A) = \text{rank}[A|B] = 3$.
    This requires the third row pivot to be non-zero:
    \[
    \lambda - 5 \neq 0 \implies \mathbf{\lambda \neq 5}, \quad \mathbf{\mu \in \mathbb{R} \text{ (any real value)}}
    \]
    \item \textbf{No Solution (Inconsistent):}
    Requires $\text{rank}(A) = 2$ and $\text{rank}[A|B] = 3$.
    This happens when $\lambda - 5 = 0$ but $\mu - 9 \neq 0$:
    \[
    \mathbf{\lambda = 5 \quad \text{and} \quad \mu \neq 9}
    \]
    \item \textbf{Infinitely Many Solutions:}
    Requires $\text{rank}(A) = \text{rank}[A|B] = 2 < 3$.
    This occurs when the entire third row is zero:
    \[
    \mathbf{\lambda = 5 \quad \text{and} \quad \mu = 9}
    \]
\end{enumerate}

\begin{markingrubric}
\textbf{Mark Distribution:}
$\bullet$ Augmented matrix and row operation: \textbf{3 Marks}\\
$\bullet$ Unique solution condition ($\lambda \neq 5, \mu \in \mathbb{R}$): \textbf{2.5 Marks}\\
$\bullet$ No solution condition ($\lambda = 5, \mu \neq 9$): \textbf{2 Marks}\\
$\bullet$ Infinite solution condition ($\lambda = 5, \mu = 9$): \textbf{2.5 Marks}
\end{markingrubric}
\end{solutionbox}

% Solution 32
\begin{solutionbox}[Solution to Question 32 (10 Marks)]
\textbf{Step 1: Cauchy-Euler Transformation}
Let $x = e^z \implies z = \ln x$, with $\theta = \frac{d}{dz}$.
Then $x y' = \theta y$ and $x^2 y'' = \theta(\theta - 1)y$.
Substituting into the equation:
\[
[\theta(\theta - 1) - 3\theta + 4]y = 2(e^z)^2 \implies (\theta^2 - 4\theta + 4)y = 2e^{2z}
\]
\textbf{Step 2: Complementary Function}
Auxiliary equation: $m^2 - 4m + 4 = 0 \implies (m - 2)^2 = 0 \implies m = 2, 2$ (repeated root).
\[
y_c = (c_1 + c_2 z)e^{2z}
\]
\textbf{Step 3: Particular Integral (Failure of Order 2)}
\[
y_p = \frac{1}{(\theta - 2)^2} (2e^{2z}) = 2 \frac{z^2}{2!} e^{2z} = z^2 e^{2z}
\]
\textbf{Step 4: Transformed Solution in $z$}
\[
y(z) = (c_1 + c_2 z)e^{2z} + z^2 e^{2z}
\]
\textbf{Step 5: Transform back to $x$}
Substitute $e^{2z} = x^2$ and $z = \ln x$:
\[
\mathbf{y(x) = (c_1 + c_2 \ln x)x^2 + x^2 (\ln x)^2}
\]

\begin{markingrubric}
\textbf{Mark Distribution:}
$\bullet$ Cauchy-Euler transformation to operator $\theta$: \textbf{3 Marks}\\
$\bullet$ Complementary function $y_c = (c_1 + c_2 z)e^{2z}$: \textbf{2.5 Marks}\\
$\bullet$ Particular integral $y_p = z^2 e^{2z}$: \textbf{2.5 Marks}\\
$\bullet$ Final solution in terms of $x$: \textbf{2 Marks}
\end{markingrubric}
\end{solutionbox}

% Solution 33
\begin{solutionbox}[Solution to Question 33 (10 Marks)]
\textbf{Part (a): Alternating Series Divergence (5 Marks)}
Given series: $\sum_{n=1}^\infty (-1)^{n-1} u_n$ where $u_n = \frac{n}{2n - 1}$.
Compute the $n$-th term limit:
\[
\lim_{n\to\infty} u_n = \lim_{n\to\infty} \frac{n}{2n - 1} = \frac{1}{2} \neq 0
\]
By the $n$-th term test for divergence: if $\lim_{n\to\infty} a_n \neq 0$, the series **MUST DIVERGE**.
Since $\lim u_n = 1/2 \neq 0$, the alternating terms oscillate between values near $+1/2$ and $-1/2$ and never settle to 0.
Thus, the series is \textbf{DIVERGENT} (fails the fundamental Leibniz condition $\lim u_n = 0$).

\textbf{Part (b): Radius and Interval of Convergence (5 Marks)}
Power series: $\sum_{n=1}^\infty \frac{(x+1)^n}{n^2 4^n}$. Here $c_n = \frac{1}{n^2 4^n}$.
\[
\frac{1}{R} = \lim_{n\to\infty} \left|\frac{c_{n+1}}{c_n}\right| = \lim_{n\to\infty} \frac{n^2 4^n}{(n+1)^2 4^{n+1}} = \frac{1}{4} \implies \mathbf{R = 4}
\]
Center of convergence is $a = -1$.
Radius $R = 4 \implies |x + 1| < 4 \implies -4 < x + 1 < 4 \implies -5 < x < 3$.
\textbf{Testing Boundaries:}
\begin{itemize}[leftmargin=6mm]
    \item At $x = 3$: $\sum_{n=1}^\infty \frac{4^n}{n^2 4^n} = \sum_{n=1}^\infty \frac{1}{n^2}$ ($p$-series with $p=2 > 1$). Converges!
    \item At $x = -5$: $\sum_{n=1}^\infty \frac{(-4)^n}{n^2 4^n} = \sum_{n=1}^\infty \frac{(-1)^n}{n^2}$. Converges absolutely by $p$-series test!
\end{itemize}
Therefore, both endpoints are included. The interval of convergence is $\mathbf{[-5, 3]}$.

\begin{markingrubric}
\textbf{Mark Distribution:}
$\bullet$ Part (a) $n$-th term test showing $\lim \neq 0$ (Divergent): \textbf{5 Marks}\\
$\bullet$ Part (b) Radius $R = 4$ and endpoint checks giving $[-5, 3]$: \textbf{5 Marks}
\end{markingrubric}
\end{solutionbox}

% Solution 34
\begin{solutionbox}[Solution to Question 34 (10 Marks)]
\textbf{Step 1: Geometric Setup}
Let the hemisphere be $x^2 + y^2 + z^2 = R^2$ with $z \ge 0$.
The base of the box lies on the flat circular equatorial plane $z = 0$, with vertices at $(\pm x, \pm y, 0)$.
The dimensions of the rectangular box are:
\[
\text{Length} = 2x, \quad \text{Width} = 2y, \quad \text{Height} = z
\]
The volume of the box is:
\[
V(x,y,z) = (2x)(2y)(z) = 4xyz
\]
Subject to the constraint $g(x,y,z) = x^2 + y^2 + z^2 - R^2 = 0$.

\textbf{Step 2: Lagrange Multipliers}
Set $\nabla V = \lambda \nabla g$:
\begin{align}
4yz &= 2\lambda x \implies 4xyz = 2\lambda x^2 \label{eq:box1} \\
4xz &= 2\lambda y \implies 4xyz = 2\lambda y^2 \label{eq:box2} \\
4xy &= 2\lambda z \implies 4xyz = 2\lambda z^2 \label{eq:box3}
\end{align}
Equating \eqref{eq:box1}, \eqref{eq:box2}, and \eqref{eq:box3}:
\[
2\lambda x^2 = 2\lambda y^2 = 2\lambda z^2 \implies \mathbf{x^2 = y^2 = z^2}
\]
\textbf{Step 3: Solving for Optimal Dimensions}
Substitute $x^2 = y^2 = z^2$ into the hemisphere equation:
\[
x^2 + x^2 + x^2 = R^2 \implies 3x^2 = R^2 \implies x = \frac{R}{\sqrt{3}}
\]
Hence:
\[
x = \frac{R}{\sqrt{3}}, \quad y = \frac{R}{\sqrt{3}}, \quad z = \frac{R}{\sqrt{3}}
\]
The box dimensions are:
\[
\text{Length} = \frac{2R}{\sqrt{3}}, \quad \text{Width} = \frac{2R}{\sqrt{3}}, \quad \text{Height} = \frac{R}{\sqrt{3}}
\]
\textbf{Step 4: Maximum Volume}
\[
V_{\max} = 4\left(\frac{R}{\sqrt{3}}\right)\left(\frac{R}{\sqrt{3}}\right)\left(\frac{R}{\sqrt{3}}\right) = \mathbf{\frac{4R^3}{3\sqrt{3}}}
\]

\begin{markingrubric}
\textbf{Mark Distribution:}
$\bullet$ Volume function $V = 4xyz$ and hemisphere constraint: \textbf{2.5 Marks}\\
$\bullet$ Lagrange equations and deduction $x^2 = y^2 = z^2$: \textbf{3.5 Marks}\\
$\bullet$ Finding $x = y = z = R/\sqrt{3}$: \textbf{2 Marks}\\
$\bullet$ Maximum volume $4R^3 / (3\sqrt{3})$: \textbf{2 Marks}
\end{markingrubric}
\end{solutionbox}

% Solution 35
\begin{solutionbox}[Solution to Question 35 (10 Marks)]
\textbf{Step 1: Spherical Polar Coordinates}
Let $x = \rho\sin\phi\cos\theta, y = \rho\sin\phi\sin\theta, z = \rho\cos\phi$.
Then $x^2 + y^2 + z^2 = \rho^2$, and $dV = \rho^2 \sin\phi \, d\rho \, d\phi \, d\theta$.
The denominator becomes:
\[
\sqrt{1 - (x^2 + y^2 + z^2)} = \sqrt{1 - \rho^2}
\]
\textbf{Step 2: Bounds for Positive Octant of Unit Sphere}
\[
\rho \in [0, 1], \quad \phi \in \left[0, \frac{\pi}{2}\right], \quad \theta \in \left[0, \frac{\pi}{2}\right]
\]
\textbf{Step 3: Integral Evaluation}
The three integrals decouple:
\[
I = \int_0^{\pi/2} d\theta \cdot \int_0^{\pi/2} \sin\phi \, d\phi \cdot \int_0^1 \frac{\rho^2}{\sqrt{1 - \rho^2}} \, d\rho
\]
\begin{itemize}[leftmargin=6mm]
    \item Angular 1: $\int_0^{\pi/2} d\theta = \frac{\pi}{2}$
    \item Angular 2: $\int_0^{\pi/2} \sin\phi \, d\phi = [-\cos\phi]_0^{\pi/2} = 1$
    \item Radial Integral: Let $\rho = \sin t \implies d\rho = \cos t \, dt$:
    \[
    \int_0^{\pi/2} \frac{\sin^2 t}{\cos t} (\cos t \, dt) = \int_0^{\pi/2} \sin^2 t \, dt = \left[ \frac{t}{2} - \frac{\sin 2t}{4} \right]_0^{\pi/2} = \frac{\pi}{4}
    \]
\end{itemize}
Multiply all three factors:
\[
I = \left(\frac{\pi}{2}\right)(1)\left(\frac{\pi}{4}\right) = \mathbf{\frac{\pi^2}{8}}
\]

\begin{markingrubric}
\textbf{Mark Distribution:}
$\bullet$ Spherical coordinate conversion and Jacobian: \textbf{3 Marks}\\
$\bullet$ Limits for the positive octant: \textbf{2 Marks}\\
$\bullet$ Trigonometric substitution for radial integral: \textbf{3 Marks}\\
$\bullet$ Final exact answer $\pi^2 / 8$: \textbf{2 Marks}
\end{markingrubric}
\end{solutionbox}

% Solution 36
\begin{solutionbox}[Solution to Question 36 (10 Marks)]
\textbf{Step 1: Statement of Gauss Divergence Theorem}
\[
\iint_S \vec{F} \cdot \hat{n} \, dS = \iiint_V (\nabla \cdot \vec{F}) \, dV
\]
Given $\vec{F} = x^3\hat{i} + y^3\hat{j} + z^3\hat{k}$ over the sphere $x^2 + y^2 + z^2 \le a^2$.

\textbf{Step 2: Volume Integral Evaluation}
Compute $\nabla \cdot \vec{F}$:
\[
\nabla \cdot \vec{F} = \frac{\partial}{\partial x}(x^3) + \frac{\partial}{\partial y}(y^3) + \frac{\partial}{\partial z}(z^3) = 3x^2 + 3y^2 + 3z^2 = 3(x^2 + y^2 + z^2)
\]
In spherical polar coordinates: $x^2 + y^2 + z^2 = \rho^2$ and $dV = \rho^2 \sin\phi \, d\rho \, d\phi \, d\theta$:
\[
\iiint_V 3\rho^2 (\rho^2 \sin\phi \, d\rho \, d\phi \, d\theta) = 3 \int_0^{2\pi} d\theta \int_0^\pi \sin\phi \, d\phi \int_0^a \rho^4 \, d\rho
\]
\[
= 3 (2\pi)(2) \left[ \frac{\rho^5}{5} \right]_0^a = 12\pi \left(\frac{a^5}{5}\right) = \mathbf{\frac{12\pi a^5}{5}}
\]

\textbf{Step 3: Surface Integral Evaluation}
On the sphere $x^2 + y^2 + z^2 = a^2$:
The outward unit normal is $\hat{n} = \frac{\vec{r}}{a} = \frac{x\hat{i} + y\hat{j} + z\hat{k}}{a}$.
\[
\vec{F} \cdot \hat{n} = \frac{x^4 + y^4 + z^4}{a}
\]
By spherical symmetry:
\[
\iint_S x^4 \, dS = \iint_S y^4 \, dS = \iint_S z^4 \, dS \implies \iint_S (x^4 + y^4 + z^4) dS = 3 \iint_S z^4 \, dS
\]
In spherical coordinates on $S$, $\rho = a, z = a\cos\phi, dS = a^2 \sin\phi \, d\phi \, d\theta$:
\[
\iint_S z^4 \, dS = \int_0^{2\pi} d\theta \int_0^\pi (a\cos\phi)^4 (a^2 \sin\phi \, d\phi) = 2\pi a^6 \int_0^\pi \cos^4\phi \sin\phi \, d\phi
\]
Let $u = \cos\phi \implies du = -\sin\phi \, d\phi$:
\[
= 2\pi a^6 \int_{-1}^1 u^4 \, du = 2\pi a^6 \left[ \frac{2}{5} \right] = \frac{4\pi a^6}{5}
\]
Therefore:
\[
\iint_S \vec{F} \cdot \hat{n} \, dS = \frac{1}{a} \left[ 3 \left(\frac{4\pi a^6}{5}\right) \right] = \frac{1}{a} \left(\frac{12\pi a^6}{5}\right) = \mathbf{\frac{12\pi a^5}{5}}
\]
Since Volume Integral ($\frac{12\pi a^5}{5}$) = Surface Integral ($\frac{12\pi a^5}{5}$), **Gauss Divergence Theorem is completely verified!**

\begin{markingrubric}
\textbf{Mark Distribution:}
$\bullet$ Stating theorem and computing $\nabla \cdot \vec{F} = 3(x^2+y^2+z^2)$: \textbf{2 Marks}\\
$\bullet$ Evaluating volume integral in spherical coordinates: \textbf{3.5 Marks}\\
$\bullet$ Evaluating surface integral via normal $\hat{n} = \vec{r}/a$: \textbf{3.5 Marks}\\
$\bullet$ Concluding verification: \textbf{1 Mark}
\end{markingrubric}
\end{solutionbox}

\newpage
